% === Begin Label Data ===
% ID: 568
% 难度星级: 
% 标签: 
% 备注: 
% 组卷引用次数: 0
% === End  Label Data ===

\begin{problem}{2009}{G}{江苏卷}{21}{函数}
按照某学者的理论，假设一个人生产某产品单件成本为 $ a $ 元，如果他卖出该产品的单价为 $ m $ 元，则他的满意度为 $ \dfrac{m}{m+a} $；如果他买进该产品的单价为 $ n $ 元，则他的满意度为 $ \dfrac{n}{n+a} $．如果一个人对两种交易（卖出或买进）的满意度分别为 $ h_1 $ 和 $ h_2 $，则他对这两种交易的综合满意度为 $ \sqrt{h_1 h_2} $．

现假设甲生产 $ A $、$ B $ 两种产品的单件成本分别为 $ 12 $ 元和 $ 5 $ 元，乙生产 $ A $、$ B $ 两种产品的单件成本分别为 $ 3 $ 元和 $ 20 $ 元，设产品 $ A $、$ B $ 的单价分别为 $ m_A $ 元和 $ m_B $ ，甲买进 $ A $ 与卖出 $ B $ 的综合满意度为 $ h_{\text{甲}} $，乙卖出 $ A $ 与买进 $ B $ 的综合满意度为 $ h_{\text{乙}} $．

（1）求 $ h_{\text{甲}} $ 和 $ h_{\text{乙}} $ 关于 $ m_A $、$ m_B $ 的表达式；当 $ m_A = \dfrac{3}{5} m_B $ 时，求证：$ h_{\text{甲}} = h_{\text{乙}} $；

（2）设 $ m_A = \dfrac{3}{5} m_B $，当 $ m_A $、$ m_B $ 分别为多少时，甲、乙两人的综合满意度均最大？最大的综合满意度为多少？

（3）记（2）中最大的综合满意度为 $ h_0 $，试问能否适当选取 $ m_A $、$ m_B $ 的值，使得 $ h_{\text{甲}} \geq h_0 $ 和 $ h_{\text{乙}} \geq h_0 $ 同时成立，但等号不同时成立？试说明理由．
\end{problem}

\begin{answer}

\end{answer}

\begin{solutions}

\end{solutions}